\subsection{The problems in prime fields}
In the remainder of the paper, we look only at the prime fields. Let $p$ be a prime and $\Fp$ be the prime field with $p$ elements. Throughout the paper, $A\lessapprox B$ means $A\le  C (\log p)^{C} B$ for some large constant $C$ (independent of $p$). $A\approx B$ means both $A\lessapprox B$ and $B\lessapprox A$.  $A\lesssim B$ means $A\le C B$ for a large constant $C$ (independent of $p$). $A\sim B$ means both $A\lesssim B$ and $B\lesssim A$. 

\subsubsection{Furstenberg set problem}

\begin{definition}[$(s,t;n,k)$-set]
    Fix integers $1\le k<n$, and numbers $0\le s\le k$, $0\le t\le (k+1)(n-k)$. Let $0<\lam\le 1$. We say $E$ is a $(s,t;n,k;\lam)$-Furstenberg set (over $\Fp$), if $E$ is a subset of $\Fp^n$ of form
    \[ E=\bigcup_{V\in\cV}Y(V). \]
    Here, $\cV\subset A(k,\Fp^n)$ is a set of $k$-planes with $\#\cV\ge \lam p^t$; for each $V\in\cV$, $Y(V)$ is a subset of $V$ with $\#Y(V)\ge \lam p^s$. 
\end{definition}

\begin{remark}
   {\rm We remark that $s,t,n,k$ are important parameters, while $\lam$ is not.
   Since $\lambda$ is not an important parameter, from now on we will drop $\lam$ from the notation and simply call $(s,t;n,k;\lam)$-set as $(s,t;n,k)$-set, thinking of $\lam$ as a fixed constant in $(0,1]$. 
   }
\end{remark}

We give the following definition.
\begin{definition}
    We say $(s,t;n,k)$ is admissible if $1\le k<n$ are integers and $0\le s\le k$, $0\le t\le (k+1)(n-k)$.
\end{definition}

The Furstenberg set problem can be formulated as follows:

\begin{question}
Suppose $(s,t;n,k)$ is admissible.
Find the largest $\bF(s,t;n,k)$, so that 
\begin{equation}\label{deff}
    \inf_{E:\ (s,t;n,k)\textup{-set}} \#E\gtrsim_{s,t,n,k,\e} p^{\bF(s,t;n,k)-\e} 
\end{equation}
holds for any $\e>0$.
\end{question}

In the above inequality, the implicit constant may depend on many parameters, but never depends on $p$. For simplicity, we may view $s,t,n,k$ as fixed parameters and write $\gtrsim_{s,t,n,k,\e}$ as $\gtrsim_\e$ throughout the paper. The dependence on $s,t,n,k$ will be mentioned when needed. 
We are interested in the behavior as $p$ goes to infinity, so the goal is to find the index $\bF(s,t;n,k)$ for fixed $(s,t;n,k)$.

We see that the problem consists of two parts:
for fixed $(s,t;n,k)$,
\begin{enumerate}
    \item (Lower bound) show for any $(s,t;n,k)$-set $E$, we have $\#E\gtrsim_{\e} p^{\bF(s,t;n,k)-\e}$;
    \item (Upper bound) show the existence of a $(s,t;n,k)$-set $E$, such that $\#E\lesssim p^{\bF(s,t;n,k)}$.
\end{enumerate}

\medskip

Based on Ren-Wang's result \eqref{renwang}, it is reasonable to make a conjecture in $\Fp^2$.

\begin{conjecture}\label{conj2d}
    Fix $s\in(0,1],t\in (0,2]$. Then for any $(s,t;2,1)$-set $E$ over $\Fp$, we have
    \[ \#E\gtrsim_{s,t,\e} p^{\min\{s+t,\frac{3}{2}s+\frac{1}{2}t,s+1  \}-\e}, \]
    for any $\e>0$.
\end{conjecture}

\begin{remark}
{\rm
    It is crucial to use the prime field $\Fp$. Otherwise, the conjecture fails. Consider the problem in $\F_q^2$, where $q=p^2$ is the square of a prime. We choose the set of lines $\cV=\{y=ax+b: a,b\in \Fp\}$, so $\#\cV=q$ and hence $t=1$. For each line $V: y=ax+b$ in $\cV$, define $Y(V):=\{(x,ax+b):x\in\Fp\}$, so $\#Y(V)=p=q^{\frac12}$ and therefore $s=\frac12$. We can calculate $\#\big(\bigcup_{V\in\cV}Y(V)\big)=p^2=q$. However, $\min\{s+t,\frac32s+\frac12t,s+1\}=\frac54$.

    Usually, people believe that the problem in a finite field is easier than in Euclidean space. However, Conjecture \ref{conj2d} remains open, whereas the Euclidean version has been resolved.
}
\end{remark}

\bigskip


In the paper, we explicitly find the function $\bF(s,t;n,k)$. The definition of $\bF(s,t;n,k)$ is tricky, so we postpone it later in Definition \ref{deffurindex}. The following are the main results.


\begin{theorem}[Upper bound]\label{upper bound}
    Suppose $(s,t;n,k)$ is admissible. Then for any prime $p$, there exits a $(s,t;n,k)$-set $E$ over $\Fp$, such that
    \[ \#E\lesssim p^{\bF(s,t;n,k)}. \] 
\end{theorem}
